% =====================================================================
% CHAPTER 2: LITERATURE SURVEY
% =====================================================================
\chapter{Literature Survey}

A comprehensive review of existing literature and related work has been conducted to understand the current state of smart glasses technology, identify key advancements in the field, and determine the research gaps that this project will aim to address. The following subsections present a critical analysis of seven significant works that form the foundation for the proposed system design.

\section{Low-Cost Smart Glasses using ESP32}
Das, Saxena, and Rout (2021) presented a low-cost smart glass implementation using the ESP-32 microcontroller at the IEEE 2nd International Conference on Applied Electromagnetics, Signal Processing, and Communication (AESPC). Their work demonstrated the feasibility of building wearable visual assistance devices using affordable microcontroller platforms. The system employed the ESP32 for basic image capture and wireless transmission, establishing that the platform possesses sufficient computational resources for real-time image acquisition and Wi-Fi communication. However, their implementation was limited to basic image capture and display without any AI-powered analysis capabilities. The system lacked scene understanding, text extraction, or any form of natural language interaction, rendering it a passive data collection device rather than an intelligent assistive tool. The proposed project will extend this foundational approach by adding a comprehensive AI processing pipeline with both cloud and local model integration.

\section{AI-Powered Smart Glasses for the Blind}
Ramesh Babu et al. (2025) proposed AI-powered smart glasses for visually impaired users using OpenCV and Python at the IEEE Wireless Antenna and Microwave Symposium (WAMS). Their system focused on real-time object detection using traditional computer vision techniques, specifically employing Haar cascades and HOG (Histogram of Oriented Gradients) descriptors for feature extraction. While effective for basic object identification tasks such as face detection and pedestrian recognition, their approach lacked the ability to provide natural language descriptions of scenes or answer contextual questions about the visual environment. The system could identify \textit{what} objects were present but could not explain \textit{what was happening} in the scene or respond to user queries. The proposed project will address this limitation by incorporating large language models capable of generating detailed, context-aware descriptions and engaging in conversational question answering about captured scenes.

\section{Deep Learning-Based Smart Glasses}
Lin et al. (2020) developed a smart glasses application system for visually impaired people based on deep learning at the Indo-Taiwan 2nd International Conference on Computing, Analytics and Networks (Indo-Taiwan ICAN). Their work employed convolutional neural networks (CNNs) for object classification, achieving high accuracy on standard benchmarks. However, their system required a powerful GPU-equipped computer with dedicated VRAM exceeding 8 GB for real-time inference, significantly limiting its portability and accessibility. The computational requirements meant that the system could only function in controlled environments with access to high-performance hardware. The proposed hybrid cloud-local architecture will overcome this constraint by offloading heavy AI tasks to free cloud APIs while keeping the local compute requirements to a minimum, enabling operation on consumer-grade hardware.

\section{Speech Recognition for Wearable Devices}
Radford et al. (2023) introduced Whisper, a robust speech recognition system trained on 680,000 hours of multilingual and multitask supervised audio data, at the 40th International Conference on Machine Learning (ICML). The Whisper model family ranges from tiny (39 million parameters, approximately 145 MB) to large (1.5 billion parameters, approximately 3 GB), offering a spectrum of accuracy-efficiency tradeoffs suitable for different deployment scenarios. The tiny model, despite its compact size, achieves competitive word error rates on standard benchmarks such as LibriSpeech and Common Voice. The proposed project will utilize the Whisper-tiny model for voice command recognition, as it provides an optimal balance between accuracy and resource consumption for the target deployment environment. The model's ability to function with minimal GPU memory will complement the Smart Routing architecture's memory management strategy.

\section{Large Language Models for Visual Understanding}
Touvron et al. (2023) published the Llama 2 family of open foundation and fine-tuned chat models from Meta AI Research. These models demonstrated that open-source large language models could achieve performance competitive with proprietary systems such as GPT-3.5 and PaLM across a wide range of natural language understanding and generation tasks. The subsequent Llama 3.3 and Llama 4 Scout models, which the proposed project will utilize through the Groq API, represent significant improvements in both language understanding and multimodal reasoning capabilities. Llama 4 Scout 17B, in particular, introduces native multimodal processing that can accept both text and image inputs, making it ideal for scene description tasks. The availability of these models through free-tier cloud APIs fundamentally changes the economics of AI-powered wearable devices.

\section{Optical Character Recognition for Accessibility}
JaidedAI (2020) developed EasyOCR, an open-source optical character recognition library supporting over 80 languages including English, Hindi, Kannada, and other Indian regional scripts. Unlike commercial OCR solutions that require paid subscriptions or API access, EasyOCR is freely available under the Apache 2.0 license and can run entirely on local hardware with modest memory requirements of approximately 180 MB. The library uses a combination of CRAFT (Character Region Awareness for Text detection) for text region detection and deep learning-based recognition networks for character identification. The recognition pipeline supports multiple scripts simultaneously and handles curved, rotated, and partially occluded text with reasonable accuracy. The proposed project will leverage EasyOCR for the text extraction capability, enabling the smart glasses to read text from signs, books, documents, and product labels---a critical feature for visually impaired users navigating public spaces.

\section{Edge Computing for AI Applications}
Chen and Ran (2019) provided a comprehensive review of deep learning with edge computing in the Proceedings of the IEEE, examining the challenges and opportunities of deploying AI models on resource-constrained edge devices. Their survey highlighted the importance of distributing AI workloads between cloud and edge devices to achieve optimal latency (sub-second response times for real-time applications), bandwidth utilization (reducing data transmission by processing locally where possible), and privacy (keeping sensitive data on-device). They identified model partitioning and workload scheduling as key technical challenges that remain actively researched. The proposed project's Smart Routing architecture will directly implement the principles outlined in their survey, demonstrating a practical application of cloud-edge AI distribution in a wearable device context with intelligent model lifecycle management.

\section{Research Gap}
The literature review reveals that while several smart glasses projects exist in the academic and commercial domains, most suffer from one or more of the following critical limitations:

\begin{enumerate}[leftmargin=1.5cm]
\item \textbf{High Cost:} Commercial solutions are priced beyond the reach of users in developing countries, and academic prototypes often use expensive hardware platforms such as Raspberry Pi 4 with dedicated GPUs.

\item \textbf{Heavy Computational Requirements:} Existing AI-powered smart glasses require high-performance GPUs with 8--12 GB of VRAM for local inference, making them incompatible with consumer-grade hardware.

\item \textbf{Limited AI Capabilities:} Most academic prototypes are restricted to a single AI modality---either object detection, or OCR, or speech recognition---without providing an integrated multi-modal experience.

\item \textbf{No Memory Optimization:} None of the reviewed works address the challenge of running multiple AI models efficiently on resource-constrained hardware through intelligent model lifecycle management.
\end{enumerate}

This project will address all four limitations through its affordable hardware design based on the ESP32-CAM, hybrid cloud-local AI architecture utilizing free-tier APIs, comprehensive multi-modal AI capabilities spanning five distinct functions, and the novel Smart Routing memory management system that will enable operation on entry-level consumer hardware.
